\begin{helper}
Fix \(\varepsilon,\eta,\varepsilon_1,\varepsilon_2,\beta\in\mathbb R\) and
\(n_0\in\mathbb N\).  Assume
\[
0<\varepsilon,\qquad 0<\eta,\qquad \eta<\varepsilon,\qquad
\beta=\frac12,
\]
and \(\noderef{ParameterHierarchy}(\eta,\varepsilon_1,\varepsilon_2,n_0)\).
Suppose that for every \(\delta>0\) there is \(N\ge n_0\) such that, for all
\(n>N\), all \(m,p\) with
\[
m=\noderef{TwoBiteNaturalReducedVertexCount}(n),\qquad
p=\noderef{TwoBiteEdgeProbability}\!\left(\frac12,n\right),
\]
and all \(I\subseteq\mathrm{Fin}(n)\) with
\[
|I|=\noderef{TwoBiteNaturalIndependenceScale}(1+\eta,n),
\]
one has
\[
\noderef{TwoBiteEventProbability}\!\left(n,m,p,\,
\{\,\omega:
I\text{ is independent in }\noderef{TwoBiteFinalGraph}(\omega)_0
\text{ and }
\noderef{TwoBiteRegularityEvent}
(k=|I|,\eta,\varepsilon_1,\varepsilon_2,\omega)\,\}\right)
\le
\delta\binom n{|I|}^{-1}.
\]
Then the probability of the event that the desired independence-number bound
fails while the \(\eta\)-scale regularity event
\[
\noderef{TwoBiteRegularityEvent}
\bigl(k=\noderef{TwoBiteNaturalIndependenceScale}(1+\eta,n),
\eta,\varepsilon_1,\varepsilon_2,\omega\bigr)
\]
holds is negligible under the two-bite law with parameters
\[
n,\quad \noderef{TwoBiteNaturalReducedVertexCount}(n),\quad
\noderef{TwoBiteEdgeProbability}(\beta,n).
\]
\end{helper}
\begin{proof}
Write
\[
M(n)=\noderef{TwoBiteNaturalReducedVertexCount}(n),\qquad
p_\beta(n)=\noderef{TwoBiteEdgeProbability}(\beta,n),\qquad
K(n)=\noderef{TwoBiteNaturalIndependenceScale}(1+\eta,n).
\]
Since \(\beta=1/2\), the law with edge parameter \(p_\beta(n)\) is the same
law as the fixed-set estimate's law with
\(p=\noderef{TwoBiteEdgeProbability}(1/2,n)\).

Apply \(\noderef{NaturalIndependenceScaleFitsTarget}\) to
\(0<\eta<\varepsilon\).  There is a threshold \(N_{\mathrm{fit}}\) such that,
for all \(n\ge N_{\mathrm{fit}}\),
\[
(K(n):\mathbb R)<
\noderef{TwoBiteIndependenceScale}(1+\varepsilon,n)
=(1+\varepsilon)\sqrt{(n:\mathbb R)\log(n:\mathbb R)}. \tag{1}
\]
Increasing the threshold if necessary and using
\(\noderef{NaturalParameterApproximation}\) together with the eventual
comparisons in \(\noderef{ParameterHierarchy}\), we may also assume
\[
0<K(n)\le n. \tag{2}
\]
Indeed, the natural scale is the ceiling of the corresponding real scale,
which is eventually positive and \(o(n)\); the stated approximation result
transfers these inequalities to the integer scale.

Fix \(\rho>0\), and set \(\delta=\rho/2\).  Applying the fixed-set estimate
with this \(\delta\), choose \(N_F\ge n_0\) such that for every \(n>N_F\) and
every \(K(n)\)-element set \(I\subseteq\mathrm{Fin}(n)\),
\[
\Pr\!\left(I\text{ is independent in }\noderef{TwoBiteFinalGraph}_0
\text{ and }R_n\right)
\le
\delta\binom n{K(n)}^{-1}, \tag{3}
\]
where \(R_n\) is the regularity event with \(k=K(n)\).  This is exactly the
event in the fixed-set estimate, because \(|I|=K(n)\).

Let \(B_n\) be the event that
\(\noderef{IndependenceNumberLess}\) at threshold
\((1+\varepsilon)\sqrt{(n:\mathbb R)\log(n:\mathbb R)}\) fails.  Unfolding
\(\noderef{IndependenceNumberLess}\), every \(\omega\in B_n\) has a finite set
\(J\subseteq\mathrm{Fin}(n)\) that is independent in the final graph and
satisfies
\[
(1+\varepsilon)\sqrt{(n:\mathbb R)\log(n:\mathbb R)}
\le (|J|:\mathbb R). \tag{4}
\]
For \(n\ge N_{\mathrm{fit}}\), (1) and (4) imply \(K(n)\le |J|\).  With
\(0<K(n)\) from (2), choose a \(K(n)\)-element subset \(I\subseteq J\).
Independence is inherited by subsets, so \(I\) is independent in the same
final graph.  Hence, for all sufficiently large \(n\),
\[
B_n\cap R_n
\subseteq
\bigcup_{\substack{I\subseteq\mathrm{Fin}(n)\\ |I|=K(n)}}
\left(\{I\text{ is independent in }\noderef{TwoBiteFinalGraph}_0\}
\cap R_n\right). \tag{5}
\]

The number of \(K(n)\)-element subsets of \(\mathrm{Fin}(n)\) is
\(\binom n{K(n)}\), and this binomial coefficient is positive by (2).  The
finite union bound applied to (5), followed by (3) for each set \(I\), gives
\[
\Pr(B_n\cap R_n)
\le
\binom n{K(n)}\,\delta\binom n{K(n)}^{-1}
=\delta
<\rho
\]
for all sufficiently large \(n\).  Since \(\rho>0\) was arbitrary, the
probability of \(B_n\cap R_n\) tends to \(0\), which is the asserted
\(\noderef{NegligibleProbability}\) statement.
\end{proof}
